% Authorship: Claude-drafted, author-revised
\section{Weights Can Emerge from Practice}

Part IV took the weights as given and asked the aggregation Berman
describes to do the work the theory needs. It does. But that leaves the
theory's central commitment untested. On his own account Berman does not say that a philosopher supplies the
weights. He says practice confers them.\footnote{Berman, \textit{Practices}, \textit{supra} note 7, at
329.} On stipulated weights alone, a model shows only how the law would come out if
a theorist set its parameters, and principled positivism makes no such claim. So the prior question is whether a weight
structure of the sort Part IV ran on can arise from practice rather than
being written down in advance.

Through the simplest of the four practice mechanisms, this Part shows that it can,
and that the weights it yields clear the same bar the stipulated ones did. That is a proof of concept, not the whole story. As implemented here the mechanism is a spare one, and Part VI shows that
richer mechanisms confer weight differently. But the possibility must be
established first, and it is established here.

The mechanism is the \textit{uptake} model of Part III. In nine-judge panels, a bench of roughly one hundred simulated judges
decides a stream of cases. Each judge on the bench carries one of four interpretive
methodologies, and the methodology sets the larger part of how readily that
judge reaches for each principle. Nothing about
the weights is stipulated. The emergent quantity is each principle's \textit{uptake share}: how often the bench
as a whole invokes it, scaled so that the most-invoked principle sits at
$1.0$. At issue is the structure that quantity takes on once the simulation has
run its two thousand cases.

Three results follow. A stable weight structure emerges, and it splits the principles into two
groups. That split does not come from any survival contest among judges,
and so answers a caution Berman has pressed. And the emergent weights meet
all four requirements, on every way of turning the bench over that I tried,
just as the stipulated weights did. A closing section marks the limits of the
mechanism, and those limits open onto Part VI.

\subsection{A Stable Two-Group Structure Emerges}

Run the bench and let it decide cases, and the uptake shares settle into
a stable arrangement. By the end, five principles stay \textit{live}, accumulating high uptake:
text, original meaning, precedent, structure, and
federalism. Those are the textual and structural principles, keyed to the
constitutional text and its architecture. Four go \textit{dead}: dignity,
equality, liberty, and welfare, the moral principles. They are not
unmentioned. Their raw uptake shares run from
about $0.35$ to about $0.50$ against $0.79$ to $1.00$ for the other five. Beyond
that difference, the live/dead cut zeroes any principle whose share falls below the
median of the nine, and the moral four fall on the wrong side of it, a point
Part VI returns to.

This membership is not a close call that tips one way
on some runs and the other way on others. Each seed is one independent run of
the simulation, and on every seed the same five are live and the same four are
dead; as the next Section develops, the same holds under every rule for
retiring and replacing judges. The structure is an output of the model, not an input to
it, and it is the same output every time. Whether a real bench would sort
the same way is not something this design can show.

Why these five? The answer is the practice-grounding idea doing its
work. The simulated bench is composed in equal quarters of four methodologies:
textualist, originalist, purposivist, and pragmatist. Each sorts the nine
principles into a band it reaches for readily, a band it reaches for
reluctantly, and a middle band. Two of them, the
textualist and the originalist, reach for the textual and structural
principles far more often than for the moral ones. Together those two methodologies make up half the bench, and both reach for
all four moral principles less readily than for anything else on the list,
while no moral principle is a favorite of more than one of the four. So the
moral principles accumulate the fewest invocations.

The principles that accumulate the most, though, are not the ones any
methodology favors most intensely but the ones the bench as a whole is willing
to take up. Although only one methodology favors precedent, it takes the
highest share of the nine, because none disfavors it; and structure survives
comfortably although none favors it at all, for the same reason. Favored
intensely by two methodologies and disfavored by the other two, original
meaning sits at the bottom of the live group. Breadth of tolerance governs,
not depth of preference.\footnote{On the static bench, averaged across the thirty seeds,
the shares are precedent $1.00$, text $0.96$, structure $0.85$, federalism
$0.82$, original meaning $0.79$, then welfare and liberty at $0.50$ and
dignity and equality at $0.35$. Because shares are scaled to the most-invoked
principle, precedent's $1.00$ records which principle leads rather than by how
much; it leads on every one of the thirty seeds, and original meaning is the
lowest live principle on twenty-eight of them. The ordering is the one the
bench-wide average of the four methodologies' invocation propensities
predicts, and that average is fixed before any case is decided: no methodology
places precedent or structure in the band it reaches for least, while each
moral principle falls into that band for at least two.}

The live/dead line falls exactly where the
composition of the bench places it. Practice confers the weight here, and argument does
not: the moral principles
are not ranked lower because they matter less, but because this bench
takes them up less. Change the composition enough and the line moves with it.
Hold everything else fixed and reseat the bench so that the purposivist
methodology, the one most disposed toward the moral principles, holds a majority
of the seats rather than a quarter, and dignity, equality, and liberty come up
live while text, original meaning, and federalism go dead.\footnote{Thirty
seeds per composition on the static bench, with the live/dead cut in place and
everything else as in the headline run. Equal quarters reproduce the headline
result on all thirty seeds. Raising the textualist share to $55\%$ leaves the
live set unchanged on all thirty and widens the cluster gap defined just below
to $0.585$, the principles a textualist bench favors having been live already.
Raising the purposivist share to $55\%$ instead trades three principles for
three on all thirty, at a cluster gap of $0.444$. Precedent and structure stay live throughout and welfare stays
dead, which is what the breadth-of-tolerance reading predicts. What
composition moves is the membership of the two groups, not the two-group
shape.} That result is worth holding onto, because it ties the emergent
structure to how a bench is staffed: who sits on the bench decides which
principles carry weight, a thread the Conclusion picks back up.

The two groups are far apart. The \textit{cluster gap} is the distance between
the average live share and the average dead share, and it runs to $0.461$ on
the bench
where no judge is ever replaced, which is the bench carrying the headline
result, and to $0.512$ averaged across all four rules for turning the bench
over.\footnote{The gap is the difference between the mean uptake share of the
live principles and the mean uptake share of the dead principles, averaged
across the thirty seeds in each of the four turnover regimes, so across one
hundred twenty runs in all; uptake shares run from $0$ to $1$, so a gap of
$0.512$ separates the two groups by more than half the available range. The
figure is pulled up by the fitness regime, which replaces the least decisive
judges with copies of the most decisive and which the design is built to do
without: on the static bench, which carries the headline result, the gap is
$0.461$ (Section V.B). The membership
of the two groups, which five principles are live, is identical on every
seed.}

A word on how firm that two-group reading is, because the honest answer has
two sides. Run by run the evidence is thinner, and two tests of it disagree
about how much thinner. Within a single run the distribution under test holds
nine numbers, one uptake share per principle, which is a thin basis for
telling one hump from two. Hartigan's dip test, which measures the largest
distance between the observed distribution and the closest single-humped one and
assumes nothing about the shape of the humps, rejects the one-hump null on
sixty-six of one hundred twenty seed-runs. A Gaussian-mixture criterion instead compares the
fit of two bell curves against one, and rejects on three.\footnote{J.A.\ Hartigan \& P.M.\ Hartigan, \textit{The Dip Test of
Unimodality}, 13 \textsc{Annals Stat.}\ 70 (1985). The mixture criterion
counts a run as two-humped when the two-component fit beats the one-component
fit by a Bayesian-information-criterion margin above ten and the component
means sit at least one standard deviation apart. The dip figure is pulled up
by the fitness regime, where it rejects on all thirty runs at a median $p$ of
$0.007$. Across the other ninety it rejects on thirty-six, and on the static
bench that carries the headline result it rejects on eleven of thirty, at a
median $p$ of $0.058$, which falls on the other side of the conventional
line.} The distance between sixty-six and three is a small sample acting
on two tests of different sensitivity, not a fact about the weights. Neither per-run
figure should carry the claim.

Pool the shares across seeds, though, and both tests agree, decisively. On
the pooled distribution the dip statistic is $0.074$ and the one-hump null
is rejected beyond any conventional level, the two-component fit beats the
one-component fit by a Bayesian-information-criterion difference of $664$
where a difference above ten is already strong, and the bimodality
coefficient, a standard index of two-groupedness, sits well into the range
that marks a genuinely split distribution. On the static bench alone, which
carries the headline result, the dip statistic rises to $0.128$. Both tests
establish that the pooled distribution is not a single hump. That is weaker
than establishing that it has exactly two.\footnote{Neither test counts
modes. The dip test rejects unimodality without saying how many humps replace
it, and the mixture criterion here compares only a two-component fit against a
one-component one. Allowed more components on the same pooled shares, the
Bayesian information criterion keeps improving past two. The two-group reading
is carried by the live/dead structure the next Sections trace, not by the
mixture comparison. The bimodality coefficient runs from $0$ to $1$, with
values above about $0.55$ taken to indicate two groups; on the pooled shares
it is $0.657$.} The claim, then, is modest: a stable two-group separation in the emergent
weights, visible in the pooled behavior of the bench and not established by
any single run. That is the claim the rest of the Part uses, and no more.

\subsection{The Separation Is Composition, Not Selection}

Berman has cautioned against reading his account through a Darwinian
lens, in which weights are fixed by some survival-of-the-fittest contest
among principles or among the judges who carry them. The caution is well
taken, and this mechanism is built to respect it, as a single comparison
shows.

The design lets the bench turn over in four ways. In
the first, the \textit{static} bench, no judge is ever replaced and none
is selected for anything: the same hundred judges decide every case from
first to last. The other three retire and replace judges over time: by appointment, reseating
a random fifth of the bench with judges drawn afresh from all four
methodologies in equal proportion; by drift, reseating a random fifth with
near-copies of other sitting judges; and by a fitness rule that retires the
least decisive judges, measuring decisiveness by the absolute size of a judge's
own force in a case, and refills their seats with copies of the most decisive.
Appointment here names one of these rules for turning this bench over, not the
appointment mechanism of Part III, which grounds weight in a different
quantity altogether. Of the four, only the fitness rule rewards success, and
it is the one whose framing Berman questions.

On the evidence here, the separation does not need selection among judges at
all. On the static bench, where no judge is ever replaced and none is selected
for anything, the two groups pull apart by a cluster gap of $0.461$. Turning the bench over by drift, or by appointment, barely moves that figure
($0.452$ and $0.472$). The fitness rule does move it, widening the gap to
$0.663$, and the reason is instructive. It does not sharpen a fixed bench. It
destroys the bench's composition: replacing the least decisive fifth with
copies of the most decisive drives the purposivist methodology extinct on
every one of the thirty seeds and the pragmatist on twenty-nine, leaving a
bench roughly four-fifths originalist and one-fifth textualist.\footnote{Final
composition after two thousand cases, thirty seeds. The static, appointment,
and drift regimes all finish within a seat or so of the twenty-five each they
began with; the fitness regime finishes at $79.7$ originalist, $20.2$
textualist, and none of the other two.} So the widest gap of the four belongs
to the bench with the least methodological diversity, exactly as the
composition comparison above would predict. The fitness regime is not evidence that
selection produces the separation. It is evidence that selection produces a
narrower bench, and a narrower bench sorts more sharply.\footnote{Cluster gaps by turnover regime, averaged across seeds: static $0.461$, drift
$0.452$, appointment $0.472$, fitness $0.663$. Tested against the static
bench, only the fitness regime differs to a degree unlikely to be chance ($p =
2\times10^{-9}$); appointment is marginal ($p = 0.043$) and drift is
indistinguishable ($p = 0.18$). The membership of the two groups is the same
under all four regimes.}

Where the split appears settles the question. If it were an artifact of a
survival contest among judges, it would show up only under the fitness rule
and be absent from the static bench, where no such contest is ever run. On the
static bench it is present at a gap of $0.461$, and the one regime
that runs a contest produces its larger gap by narrowing the bench rather than
by ranking principles. The separation traces to the composition of the bench, to the mix of
methodologies and how readily each reaches for each principle, rather than to
any winnowing of judges as such. Winnowing matters only where it changes that
mix. The fitness regime changes it; appointment and drift do not. So
the live/dead line is drawn by which principles a given bench is disposed to
take up, a standing fact about the practice rather than the outcome of a
contest. This is arguably the reading Berman's caution asks for. Nothing in the result depends on selecting among
judges, and the headline separation is the one that holds on the bench where
no judge is ever retired.

One reinforcement does remain, and it should be named rather than left in the
code. A judge's readiness to invoke a principle is not fixed by methodology
alone: a small part of it tracks the weights the bench currently carries, so a
principle the live/dead cut has zeroed is thereafter reached for less often.
That is a rich-get-richer dynamic among principles, and it runs on the static
bench along with the others. Its effect is measurable and modest. The five
live principles settle exactly where the methodology mix alone would put them,
while the four dead ones settle about a fifth lower than it
would.\footnote{Comparing the emergent shares against the bench-wide average
invocation propensity, which is fixed before any case is decided, the five
live principles come in at $1.00$ of the predicted value on average and the
four dead ones between $0.795$ and $0.798$.} So
the contest among judges that Berman's caution rules out is absent. In its
place is a feedback among principles that widens a gap the composition of the
bench had already created.

\subsection{The Emergent Weights Meet the Four Requirements}

Practice-grounded weights must clear the same bar as stipulated ones, and,
held to the four requirements of Part III, they seem to.
On the runs reported here they meet all four, and under every one of the four
turnover regimes.

First, they are \textit{well-defined}: the surviving set is nondegenerate and
identical on every seed and under every turnover regime. How many survive,
five here, is not itself informative, since the live/dead cut sits at the
median of nine and so leaves five standing as a matter of arithmetic. The
requirement asks only that at least two principles carry weight above a small
cutoff, and what these runs add beyond that is the stability of the surviving
set rather than its size.

Second, they are \textit{noise-tolerant}. The emergent weights absorb Gaussian
weight error up to a critical uncertainty, the level at which more than one
case in five changes its verdict, and that level averages between $0.091$ and
$0.097$ depending on the turnover regime, which sits above the $0.07$ floor of
Part III.\footnote{Those two figures are the spread of the four regime means. The
requirement is scored run by run, and run by run the emergent critical
uncertainty ranges from $0.075$ to $0.113$ across the hundred and twenty runs,
so the worst single run clears the floor by $0.005$ rather than by the margin
the regime means suggest. Recall too that the floor is the lowest cell mean of
Part III's forty-five-configuration grid, not the stipulated engine's worst
run, so clearing it does not establish that an emergent mechanism is more
noise-tolerant than the stipulated engine at its worst.} Third, they
\textit{weigh magnitude} rather than merely rank, passing the magnitude probe,
which perturbs only the magnitudes of the activations and asks whether the
verdict changes; the stipulated engine passed the same probe. Fourth, they deliver
\textit{more law than the hierarchical Hart}: the determinacy margin over
that charitable baseline runs to about $0.267$ to $0.271$, positive on
every seed, the same $+0.27$ the stipulated engine posted in
Part IV.\footnote{The requirements pass on every seed under all four
turnover regimes: well-defined, noise-tolerant, magnitude-weighing, and
more-law-than-Hart each hold for the static, drift, appointment, and fitness
benches alike. The margin over the hierarchical Hart is
$0.267$ to $0.271$ across regimes, matching the stipulated engine's
$+0.272$ from Section IV.B.}

The practice-grounded margin lands on the stipulated engine's own figure, but
that coincidence carries less information than it appears to. The hierarchical
Hart consults only the activations and the category tiers, never the weights,
so its determinacy rate is the same whichever weight vector it is compared
against. The margin is therefore the weighted sum's own determinacy rate less
a constant, and that rate is nearly flat across any weight vector that is not
degenerate. The two margins could hardly have failed to agree here. The
requirement establishes that the emergent weights are nondegenerate. That is
worth knowing, though the first requirement had already asked nearly the same
thing. The
margin is also measured on independent activations alone, and Part IV reports
that it reverses under the correlated ones.

Across the full weight vector the spread is wide. The vector's coefficient of
variation, a scale-free measure of how spread out the weights are, sits near $0.90$. But that spread is the live/dead split showing up as a
statistic, not graded variation among the live principles.\footnote{Among the
five that survive, the weights run from $0.79$ to $1.00$ and their coefficient
of variation is $0.09$, which is tighter than any of the other three
mechanisms manages across its own nine. \textit{See infra} Part~VI.B.}

This matters for the scope of the claim. Part III
was explicit that gradation, weights finely spread rather than
clustered, is measured but not required, because Berman permits
principles to carry equal weight.\footnote{Berman, \textit{Practices},
\textit{supra} note 7, at 323--24.} So a structure that sorts principles
into two groups without grading finely within them is not a defect to be
explained away. It meets every requirement the theory imposes, and those requirements are
arguably the standard. Whether the two-group shape is the only
shape practice can produce is a separate question, and its answer is no:
Part VI's other mechanisms grade their weights, and confer them
differently, which is the point of running more than one.

Two further features of the emergent weights, neither a requirement, are
worth recording. The first locates what drives the separation: how far the two
groups pull apart tracks how differently the four
methodologies reach for principles. Compress those differences, make the
methodologies reach for principles at more nearly equal rates, and the cluster gap narrows. Read off the raw uptake shares with the live/dead cut
switched off, the gap runs $0.135$ when the methodologies are close, $0.257$
when moderately spread, and $0.361$ at the wide spread the headline bench
uses. The size of the separation, then, is driven by how far apart the
methodologies' invocation tendencies lie, a different lever from the mix of
seats the composition sweep varied.\footnote{Cluster gap by the spread of
invocation tendencies across methodologies, measured with the live/dead cut
switched off: narrow $0.135$, medium $0.257$, wide $0.361$. These sit below the $0.461$ static-bench figure reported above, and the reason
is not that the cut zeroes anything: both figures are read off the uptake
shares before any principle is set to zero. What separates them is that the
cut, where it is switched on, operates while the simulation runs, so a
principle it zeroes is invoked less often in the cases that follow and the two
groups draw further apart than they otherwise would. Where the cut is placed
then governs how many principles count as live, seven surviving a lenient
placement, five the default, three a stringent one, but the underlying gap
stays near $0.44$ to $0.46$ throughout, so the clean live/dead binary is in
part an effect of where the cut falls, not a fact about the weights alone. How often the
bench turns over barely matters (a gap near $0.37$ whether it turns over
every $250$, $500$, or $1{,}000$ cases).}

Second, the bench co-invokes within categories. The distinction from Part IV
matters and is easy to lose. Part IV's correlated cases are cases in which the
activations move together, and the uptake bench does not generate those: its
case stream is drawn independently, exactly as Part IV's conservative baseline
was. The bench generates instead a pattern in which judges reach for
principles together, a fact about invocation rather than about the cases.

Nor does that pattern track the three categories evenly. Within a category,
principles are co-invoked at $0.80$ for the textual pair, $0.64$ for the
structural, and $0.55$ for the moral, on a scale where $1.0$ would mean that a
judge reaching for either always reaches for the other. Across categories the figure depends
entirely on which pairing is taken: textual with moral is $0.21$, the most
separated pairing the methodology table can produce, but textual with
structural is $0.66$, which is higher than the moral principles manage among
themselves, and structural with moral is $0.38$. Averaged over all three
cross-category pairings the ratio of within to across is about $1.6$, not the
three to one the most separated pairing suggests.\footnote{Ten seeds on the
static bench. The structure is two blocs rather than three categories: the
textual and structural principles are co-invoked at nearly the within-category
rate, and both are separated from the moral principles. That is the same
division as the live/dead line, which is what one would expect if both are
consequences of the same methodology mix.} Since the tendency to co-invoke within a category is
partly built into the methodology table the bench starts from, this is a
consistency check rather than an independent result.

Beyond those two, one measurement remains. It records how far
the panel's aggregate behavior sits from the weighted sum it is built from. It is not a test of whether the bench runs
a weighted sum, and should not be read as one: each judge's force is
constructed as a weighted sum of the principles it invokes, using the emergent
weights, so there is no rival rule it could have been running and a high
correlation was guaranteed by construction.

The comparison measures instead how far the panel's force falls short of the
full weighted sum. Because each judge sums only the principles it happens to invoke, the
panel's mean force is a thinned version of the full weighted sum over all
nine, and it runs about two-thirds the size. Correlation between the two is
$0.970$ and the root-mean-square difference is $0.42$.\footnote{The shrinkage
comes from the invocation rate and not from the determinacy threshold, which enters nowhere in
either quantity compared and which in any event could not alter a force, only
whether a case is called. Over the live principles the mean invocation rate is
about $0.63$, and the panel force runs at about that fraction of the full
weighted sum.} The one figure here that could have come out otherwise is
verdict agreement: the panel and the full weighted sum return the same verdict
on $85.9\%$ of cases, so invoking a subset rather than the whole list changes
the outcome about one case in seven. That is the quantity worth carrying
forward.

\subsection{What the Mechanism Does Not Yet Reach}

Emergence is therefore possible: a weight structure of the kind Part IV
ran on can arise from a simulated practice, meet every requirement, and match
the stipulated engine's Hart margin, without a theorist setting a
single weight, a match this Part has already said carries less information
than it appears to. That is the proof of concept this Part was asked to supply, though a proof
of concept is all it should be taken for.

Its limits mark the way to Part VI. The quantity that emerged here is
uptake share, how often the bench reaches for a principle, which is the
sparest reading of Berman's ``taking-up.''\footnote{Berman,
\textit{Practices}, \textit{supra} note 7, at 329.} In his fuller sense it is not yet weight, which a principle earns by
prevailing in contested disputes and by embedding itself in the holdings
that go on relying on it. Bare frequency of invocation is presumably a thinner thing than conferred
authority, and one would not expect the two to coincide. The reliance mechanism of Part VI models that fuller notion, and the remaining two,
the strategic and appointment mechanisms, model still others, none of which would likely
reduce to invocation counts. So this Part settles the question of possibility with the simplest
mechanism available. It leaves open the question that Part VI makes
central, whether the several mechanisms of
practice, asked to confer weight on the same principles, confer the same
weight. They do not.
